% Libro 1 · anexos: formulario, plan de estudio y simulacro final
\cleardoublepage
\makeatletter\renewcommand{\zuka@mundoactual}{Formulario}
\phantomsection\addcontentsline{toc}{chapter}{Formulario}
{\sffamily\fontsize{26}{30}\selectfont\firalight Formulario\par}\vspace{2pt}
{\color{zlime}\rule{18mm}{1.6pt}}\par\vspace{6pt}
{\small\color{zmute}Todo lo que hay que tener en la cabeza para este libro, en dos páginas. Si una fórmula de aquí no te
dice nada, vuelve a la lección indicada.\par}

\begin{multicols}{2}\small
\kicker{Números (1, 6)}
\begin{itemize}
  \item $a - (-b) = a + b$. Signos: par de menos $\to +$.
  \item Jerarquía: ( ) $\to$ potencias $\to$ $\times\div$ $\to$ $+-$.
  \item $\tfrac ab + \tfrac cd = \tfrac{ad + bc}{bd}$; $\tfrac ab \div \tfrac cd = \tfrac{ad}{bc}$.
  \item $0.\overline{ab} = \tfrac{ab}{99}$.
  \item $a^m a^n = a^{m+n}$, $(a^m)^n = a^{mn}$, $a^{-n} = \tfrac{1}{a^n}$.
  \item Divisores de $p^a q^b$: $(a + 1)(b + 1)$.
  \item $\text{MCD}\cdot\text{MCM} = a \cdot b$.
  \item Últimas cifras: 2 (2486), 3 (3971), 7 (7931), 8 (8426), 4 (46), 9 (91).
\end{itemize}

\kicker{Aritmética (2)}
\begin{itemize}
  \item Suma $=$ promedio $\times$ cantidad.
  \item Promedio conjunto $= \tfrac{n_1 \bar x_1 + n_2 \bar x_2}{n_1 + n_2}$.
  \item Peor caso: $n(k - 1) + 1$ para asegurar $k$ iguales entre $n$ tipos.
  \item $\overline{ab} = 10a + b$; $\overline{ab} + \overline{ba} = 11(a + b)$.
\end{itemize}

\kicker{Proporciones (4)}
\begin{itemize}
  \item Razón $a : b$, total $N$: parte $= \tfrac{N}{a + b}$.
  \item Directa: cociente fijo. Inversa: producto fijo.
  \item Juntos: $\tfrac{ab}{a + b}$. Encuentro: $\tfrac{d}{v_1 + v_2}$.
  \item $+p\,\%$: $\times(1 + \tfrac{p}{100})$; sucesivos se multiplican.
  \item Sin IVA: $\div 1.19$.
\end{itemize}

\columnbreak
\kicker{Ecuaciones (5)}
\begin{itemize}
  \item Consecutivos: $n - 1$, $n$, $n + 1$.
  \item Edades: la diferencia es constante.
  \item Suma $S$ y diferencia $D$: mayor $\tfrac{S + D}{2}$.
  \item Desigualdad $\times$ negativo: se invierte.
\end{itemize}

\kicker{Geometría (7)}
\begin{itemize}
  \item Triángulo: $180^\circ$; polígono: $(n - 2)180^\circ$.
  \item $c^2 = a^2 + b^2$; ternas 3-4-5, 5-12-13, 8-15-17, 7-24-25.
  \item $45$-$45$-$90$: $x, x, x\sqrt2$. $30$-$60$-$90$: $x, x\sqrt3, 2x$.
  \item Semejanza $k$: áreas $k^2$, volúmenes $k^3$.
  \item Triángulo $\tfrac{bh}{2}$; trapecio $\tfrac{(B + b)h}{2}$; rombo $\tfrac{Dd}{2}$.
  \item Círculo: $2\pi r$, $\pi r^2$; sector $\tfrac{\theta}{360}$.
  \item Prisma $A_b h$; cono y pirámide $\tfrac13 A_b h$; esfera $\tfrac43 \pi r^3$.
  \item 1 L $= 1\,000$ cm$^3$; 1 m$^3 = 1\,000$ L.
\end{itemize}
\end{multicols}

\clearpage\renewcommand{\zuka@mundoactual}{Plan de estudio}
\phantomsection\addcontentsline{toc}{chapter}{Plan de estudio con la app}
{\sffamily\fontsize{26}{30}\selectfont\firalight Plan de ocho semanas\par}\vspace{2pt}
{\color{zlime}\rule{18mm}{1.6pt}}\par\vspace{6pt}
{\small Este plan combina el libro con la app. Cada día: una lección del libro (30 a 40 minutos), su video y sus
juegos en Zuka (20 minutos) y el taller. Los sábados, la evaluación del mundo con reloj. Si ya dominas un mundo, haz solo
su evaluación y pasa al siguiente.\par}
\vspace{6pt}
\begin{center}\small
\begin{tabularx}{\linewidth}{@{}>{\sffamily\bfseries\color{zolive}}lX>{\raggedleft\arraybackslash}p{24mm}@{}}
\toprule
\textbf{Semana} & \textbf{Lecciones} & \textbf{Evaluación}\\\midrule
1 & 1.1 Enteros · 1.2 Fracciones · 1.3 Decimales y porcentajes · 1.4 Potencias & Mundo 1\\
2 & 2.1 Lo que queda · 2.2 Promedios · 2.3 Criptoaritmética · 2.4 Máximos y mínimos & Mundo 2\\
3 & 3.1 Operaciones nuevas · 3.2 Anidados · 4.1 Razones y repartos & Mundo 3\\
4 & 4.2 Regla de tres · 4.3 Velocidad y trabajo · 4.4 Porcentajes · 4.5 Aumentos e IVA & Mundo 4\\
5 & 5.1 Del texto a la ecuación · 5.2 Edades y mezclas · 5.3 Sistemas & Mundo 5\\
6 & 6.1 Primos · 6.2 MCM y MCD · 6.3 Residuos · 7.1 Ángulos & Mundo 6\\
7 & 7.2 Pitágoras · 7.3 Semejanza · 7.4 Áreas · 7.5 Sombreadas · 7.6 Volúmenes & Mundo 7\\
8 & Repaso de errores (los ejercicios marcados) · Simulacro final de este libro & Simulacro\\
\bottomrule
\end{tabularx}
\end{center}

\begin{tcolorbox}[enhanced,colback=zsoft,boxrule=0pt,arc=3pt,left=5mm,right=5mm]
\kicker{Cuaderno de errores}\par\vspace{2pt}
Cada vez que falles un ejercicio, anota en una libreta: la lección, qué hiciste, cuál era la idea y la
\emph{trampa} en la que caíste. Antes del examen, lee solo esa libreta. Es el repaso más rentable que existe, porque está
hecho a tu medida.
\end{tcolorbox}

\clearpage\renewcommand{\zuka@mundoactual}{Simulacro final}
\setcounter{ej}{0}
\phantomsection\addcontentsline{toc}{chapter}{Simulacro final}
\immediate\write\zuka@sol{\noexpand\solmun{Simulacro final}}
\cabecera{SIMULACRO FINAL · LIBRO 1}{Número y forma}{20 preguntas · 30 minutos\\sin calculadora}\vspace{6pt}
\begin{ejercicios}
\item $\dfrac{(-3)^2 - 3^2}{(-3)^3}$ vale
\opciones{$0$}{$\tfrac{2}{3}$}{$-\tfrac23$}{$1$}
\sol{A. El numerador es $9 - 9 = 0$. (1.1)}
\item Un tanque está lleno hasta $\tfrac{2}{3}$. Si se sacan 40 litros, queda lleno hasta $\tfrac{1}{4}$. Su capacidad es
\opciones{96 L}{120 L}{80 L}{160 L}
\sol{A. $\tfrac23 - \tfrac14 = \tfrac{5}{12}$ son 40 L. (1.2)}
\item ¿Cuál es el mayor?
\opciones{$0.\overline{3}$}{$33\,\%$}{$\tfrac{33}{100}$}{$0.333$}
\sol{A. $\tfrac13$ es mayor que $0.333$. (1.3)}
\item $\dfrac{2^{30} + 2^{30}}{2^{29}}$ es igual a
\opciones{4}{2}{$2^{31}$}{1}
\sol{A. $\tfrac{2^{31}}{2^{29}}$. (1.4)}
\item El promedio de 8 números es 15. Si se elimina el 1, el promedio de los demás es
\opciones{17}{16}{14}{15}
\sol{A. $(120 - 1)/7 = 17$. (2.2)}
\item ¿Cuántas personas debe haber para asegurar que dos nacieron el mismo día de la semana?
\opciones{7}{8}{14}{2}
\sol{B. (2.4)}
\item Si $a \triangle b = a^2 - ab$, entonces $(2 \triangle 3) \triangle 1$ vale
\opciones{6}{$-2$}{2}{0}
\sol{A. $2 \triangle 3 = -2$; $(-2) \triangle 1 = 4 + 2 = 6$. (3.1)}
\item Se reparten 180 manzanas en razón $1 : 2 : 3$. La mayor parte es
\opciones{60}{90}{120}{30}
\sol{B. (4.1)}
\item Tres pintores pintan un muro en 6 horas. ¿Cuánto tardan 9 pintores?
\opciones{18 h}{2 h}{3 h}{1 h}
\sol{B. 18 pintor-horas. (4.2)}
\item Una llave llena un tanque en 6 h y un desagüe lo vacía en 9 h. Con los dos abiertos se llena en
\opciones{18 h}{3 h}{15 h}{3.6 h}
\sol{A. $\tfrac16 - \tfrac19 = \tfrac{1}{18}$. (4.3)}
\item El $30\,\%$ de un número es 24. El $45\,\%$ de ese número es
\opciones{36}{32}{40}{54}
\sol{A. El número es 80. (4.4)}
\item Un artículo aumenta $10\,\%$ y luego se rebaja $10\,\%$. Si al final cuesta \$99\,000, costaba al principio
\opciones{\$99\,000}{\$100\,000}{\$98\,010}{\$101\,000}
\sol{B. Factor $0.99$. (4.5)}
\item La suma de dos números es 40 y su diferencia 12. El mayor es
\opciones{26}{28}{24}{14}
\sol{A. $\tfrac{40 + 12}{2}$. (5.3)}
\item Hoy Ana tiene 30 años y su hijo 6. ¿Dentro de cuántos años Ana tendrá el triple de la edad de su hijo?
\opciones{4}{6}{8}{12}
\sol{B. $30 + x = 3(6 + x)$: $x = 6$. (5.2)}
\item ¿Cuántos divisores positivos tiene 100?
\opciones{6}{8}{9}{10}
\sol{C. $2^2 5^2$: $3 \cdot 3$. (6.1)}
\item Dos señales se repiten cada 8 y 12 minutos. Si coinciden a las 9:00, vuelven a coincidir a las
\opciones{9:20}{9:24}{9:48}{9:12}
\sol{B. MCM 24. (6.2)}
\item La última cifra de $13^{2026}$ es
\opciones{3}{9}{7}{1}
\sol{B. Como la de $3^{2026}$: $2026 \bmod 4 = 2$. (6.3)}
\item Los ángulos de un triángulo miden $x$, $2x$ y $x + 20^\circ$. El mayor mide
\opciones{$40^\circ$}{$80^\circ$}{$60^\circ$}{$100^\circ$}
\sol{B. $4x + 20 = 180$, $x = 40$. (7.1)}
\item El área de un triángulo rectángulo de hipotenusa 10 y un cateto 6 es
\opciones{24}{30}{48}{60}
\sol{A. El otro cateto es 8. (7.2)}
\item Un balón de radio $r$ se reemplaza por otro de radio $2r$. Su volumen se multiplica por
\opciones{2}{4}{6}{8}
\sol{D. (7.6)}
\end{ejercicios}
\makeatother
